\chapter{Introduction}
\label{chapter:introduction}

\section{Purpose of This Document}

This document presents the Week 2 architectural foundation for the ICT Infrastructure Case Study. The assignment requires the role of an Enterprise Infrastructure Architect, responsible for researching, designing, justifying and advising on enterprise infrastructure and monitoring solutions.

The approach taken in this document is grounded in TOGAF as the enterprise architecture framework and ArchiMate as the modelling language. Rather than treating infrastructure design as a standalone technical exercise, all decisions are derived from and traced back to the business context, stakeholder needs and identified risks.

This chapter establishes the enterprise context (Step 0), which serves as the foundation for all architectural decisions in subsequent chapters.

\section{Assumptions and Scope}

All infrastructure and monitoring designs in this document are based on the enterprise context described below. Where assumptions were made, they are explicitly stated. No implementation or system administration tasks are within scope --- this document focuses on research, design and architectural advice.

% ============================================================
\chapter{Enterprise Context}
\label{chapter:enterprise-context}
% ============================================================

\section{Organisation Overview}

The organisation used throughout this case study is \textbf{LogiFlux B.V.}, a fictitious Dutch regional wholesale distribution company operating in the Fast-Moving Consumer Goods (FMCG) sector. LogiFlux acts as an intermediary between product suppliers and retail customers, managing the full chain from goods intake and warehousing through to order fulfilment and delivery.

\begin{itemize}
    \item \textbf{Legal form:} Private limited company (B.V.), incorporated in the Netherlands
    \item \textbf{Sector:} Wholesale distribution --- FMCG
    \item \textbf{Founded:} 2009
    \item \textbf{Annual revenue (estimated):} \texteuro 42 million
    \item \textbf{Customers:} Retail chains, independent supermarkets and hospitality businesses
\end{itemize}

\section{Locations}

LogiFlux operates from two permanent physical sites:

\begin{itemize}
    \item \textbf{Eindhoven (HQ + Primary Distribution Centre):} The main site houses the executive team, IT department, finance, HR, customer service and the primary warehouse and distribution operations.
    \item \textbf{Groningen (Secondary Distribution Centre):} A second warehouse and distribution facility serving customers in the northern Netherlands. This site operates under the management of a regional operations manager reporting to Eindhoven.
\end{itemize}

The two sites are approximately 220~km apart and are connected via a leased MPLS WAN link with a redundant internet-based VPN failover path.

\section{Organisational Structure}

LogiFlux employs approximately 200 staff distributed as follows:

\begin{table}[h]
\centering
\begin{tabular}{lll}
\toprule
\textbf{Department} & \textbf{Headcount} & \textbf{Primary Site} \\
\midrule
Distribution Centre (warehouse, logistics, drivers) & 120 & Both (70 EHV / 50 GRN) \\
Customer Service & 24 & Eindhoven \\
Administration (Finance \& HR) & 12 & Eindhoven \\
IT & 10 & Eindhoven \\
Management & 8 & Eindhoven \\
Maintenance \& Support & 26 & Both \\
\midrule
\textbf{Total} & \textbf{200} & \\
\bottomrule
\end{tabular}
\caption{Personnel distribution across departments and sites}
\label{tab:personnel}
\end{table}

\section{Core Business Processes}

The following core business processes drive all infrastructure and monitoring requirements:

\begin{enumerate}
    \item \textbf{Customer Order Management} --- Retail customers place orders via the customer portal (web-based) or by telephone via customer service representatives. Orders are processed, validated against inventory and scheduled for fulfilment.

    \item \textbf{Goods Storage and Distribution} --- Warehouse staff receive, store and pick goods across both distribution centres. The system manages inventory levels, triggers replenishment and coordinates outbound logistics.

    \item \textbf{Financial Administration and Accounting} --- The finance department processes invoices, manages accounts payable and receivable, and produces financial reporting. This process is subject to Dutch regulatory requirements (Belastingdienst, GDPR).

    \item \textbf{Salary Administration and Personnel Registration} --- HR manages contracts, payroll, attendance records and personnel data. This is a high-sensitivity process due to the personal and financial data involved.
\end{enumerate}

\section{Business Criticality and Risk Profile}

LogiFlux operates in a time-sensitive distribution environment. Order processing and warehouse management systems must be available during operational hours, and the customer portal must be accessible at all times.

\begin{table}[h]
\centering
\begin{tabular}{p{4.5cm}p{2cm}p{7cm}}
\toprule
\textbf{Business Process} & \textbf{Criticality} & \textbf{Risk / Impact of Downtime} \\
\midrule
Customer Order Management & High & Loss of orders, SLA breach, revenue impact, customer churn \\
Goods Storage \& Distribution & High & Operational standstill, delivery delays, contractual penalties \\
Financial Administration & Medium & Delayed invoicing, regulatory non-compliance \\
Salary \& HR Administration & Medium & Payroll errors, GDPR non-compliance, staff dissatisfaction \\
\bottomrule
\end{tabular}
\caption{Business process criticality and risk overview}
\label{tab:criticality}
\end{table}

Key risk factors for LogiFlux include:

\begin{itemize}
    \item \textbf{Seasonal peak loads:} Customer traffic can spike from approximately 1{,}200 to 10{,}000 transactions per day during promotional periods (e.g. holidays, end-of-quarter campaigns). Infrastructure must handle these peaks without degradation.
    \item \textbf{Single point of failure at Groningen:} The secondary site has limited local IT staff, making remote management and monitoring critical.
    \item \textbf{Data sensitivity:} Financial records, payroll data and customer purchase histories are all subject to GDPR and must be protected.
    \item \textbf{Supply chain dependencies:} Integration with supplier systems via EDI creates external dependency risks.
    \item \textbf{Ransomware and business email compromise:} As a mid-size distribution company with financial transaction exposure, LogiFlux is a realistic target for financially motivated attackers.
\end{itemize}

\section{Key Stakeholders and Their Interests}

\begin{table}[ht]
\centering
\begin{tabularx}{\textwidth}{lX}
\toprule
\textbf{Stakeholder} & \textbf{Interests and Concerns} \\
\midrule
CEO / Board & Business continuity, cost efficiency, regulatory compliance, reputation \\
Operations Director & Availability of order and warehouse systems during operational hours \\
IT Director & Infrastructure reliability, security posture, manageability, scalability \\
Finance Manager & Data integrity of accounting systems, audit trails, regulatory compliance \\
HR Manager & Privacy and integrity of personnel data (GDPR), payroll accuracy \\
Customer Service Team & Reliable access to the customer order system, fast issue resolution \\
Retail Customers & Order accuracy, portal availability, delivery reliability, data privacy \\
Distribution Staff & Reliable warehouse management systems, minimal downtime during shifts \\
\bottomrule
\end{tabularx}
\caption{Key stakeholders and their interests}
\label{tab:stakeholders}
\end{table}

\section{Customer Load and Capacity Assumptions}

\begin{itemize}
    \item \textbf{Average daily load:} approximately 1{,}200 customer transactions
    \item \textbf{Peak daily load:} up to 10{,}000 customer transactions (seasonal spikes)
    \item \textbf{Peak-to-average ratio:} approximately 8:1 --- infrastructure must be designed to handle peak load without significant over-provisioning during baseline periods
    \item \textbf{Operating hours:} Monday--Friday, 07:00--18:00 for warehouse operations; customer portal available 24/7
    \item \textbf{Internal users:} approximately 200 concurrent users maximum, with peak usage during morning shift handover (07:00--09:00)
\end{itemize}

\section{Summary of Assumptions}

The following assumptions underpin this entire document:

\begin{enumerate}
    \item LogiFlux B.V. is a fictional organisation. Any resemblance to real companies is coincidental.
    \item Both physical sites are in the Netherlands; all regulations referenced are Dutch/EU.
    \item The IT department (10 staff) is responsible for all infrastructure, with no dedicated security operations centre (SOC). Security monitoring is managed by the IT team.
    \item No existing enterprise architecture documentation is assumed to exist; this document represents a greenfield design.
    \item Cloud services may be used where cost-effective and appropriate; on-premises infrastructure is preferred for sensitive data.
    \item The leased MPLS link between sites provides the primary inter-site connectivity; an internet-based VPN acts as failover.
\end{enumerate}

% ============================================================
\chapter{Architecture Method}
\label{chapter:architecture-method}
% ============================================================

\section{Overview}

This document applies \textbf{TOGAF} (The Open Group Architecture Framework) as its guiding enterprise architecture method, with \textbf{ArchiMate} as the modelling language for expressing architectural views. These are not merely referenced in name --- they actively structure how analysis is conducted, how decisions are documented, and how diagrams are constructed.

\section{TOGAF and the Architecture Development Method}

TOGAF's core process is the \textbf{Architecture Development Method (ADM)}, an iterative cycle of phases for developing and maintaining enterprise architecture. For this case study, the following phases are directly applied:

\begin{table}[ht]
\centering
\begin{tabularx}{\textwidth}{lX}
\toprule
\textbf{ADM Phase} & \textbf{Application in This Document} \\
\midrule
Preliminary & Establishing the architectural approach, scope, and principles. Defines that TOGAF and ArchiMate are used and why. \\
Phase A --- Architecture Vision & Step 0 (Enterprise Context): defining the organisation, stakeholders, business goals and risk profile that drive all architectural decisions. \\
Phase B --- Business Architecture & Step 3 (Requirements Analysis): deriving requirements from business processes and stakeholder concerns. \\
Phase C --- Information Systems Architecture & Step 2 (Infrastructure Design): mapping business services to application and data components. \\
Phase D --- Technology Architecture & Step 2 (Infrastructure Design): designing the physical and virtual infrastructure, network segmentation, server clusters and connectivity. \\
Phase E --- Opportunities \& Solutions & Step 4 (Monitoring Design): identifying monitoring solutions that address gaps identified during the architecture phases. \\
\bottomrule
\end{tabularx}
\caption{TOGAF ADM phases applied in this document}
\label{tab:adm}
\end{table}

TOGAF's principle of \textbf{traceability} is applied throughout: every infrastructure component and monitoring decision is traceable back to a business requirement, stakeholder concern, or identified risk defined in the enterprise context.

\section{ArchiMate as Modelling Language}

ArchiMate is an open modelling language standardised by The Open Group, designed to describe, analyse and visualise enterprise architectures across multiple domains. It is used in this document to produce architecture diagrams that are consistent, unambiguous and layered.

ArchiMate organises elements across three \textbf{layers}:

\begin{itemize}
    \item \textbf{Business Layer} --- represents business processes, roles, functions and services as experienced by stakeholders (e.g. Customer Order Management, Salary Administration).
    \item \textbf{Application Layer} --- represents the software applications and services that support business processes (e.g. order management system, IAM/SSO, ERP).
    \item \textbf{Technology Layer} --- represents the physical and virtual infrastructure that hosts applications (e.g. server clusters, VLANs, firewalls, storage, WAN links).
\end{itemize}

Relationships between elements (such as \emph{realisation}, \emph{serving}, \emph{assignment} and \emph{composition}) are used to make dependencies and interactions explicit. This prevents the common pitfall of designing infrastructure in isolation from the business functions it must support.

\section{How the Method Influenced This Work}

The choice of TOGAF and ArchiMate had concrete effects on the structure and content of this document:

\begin{enumerate}
    \item \textbf{Enterprise context first.} Following ADM Phase A, the enterprise context was defined before any technical design began. Infrastructure choices are consequences of business needs, not the other way around.
    \item \textbf{Layered diagrams.} All architecture diagrams use ArchiMate notation to clearly separate business, application and technology concerns, making dependencies visible and reviewable.
    \item \textbf{Stakeholder-driven requirements.} ADM Phase B drives the requirements analysis --- each requirement is attributed to a stakeholder, business process or risk, rather than being a generic technical wish-list.
    \item \textbf{Justified decisions.} TOGAF's emphasis on architectural decisions with rationale means that every significant infrastructure choice in this document is accompanied by a justification and trade-off discussion.
\end{enumerate}
